\documentclass[11pt]{article}
\usepackage[margin=1in]{geometry}
\usepackage{amsmath,amssymb}
\usepackage{booktabs}
\usepackage{longtable}
\usepackage{hyperref}
\usepackage{xcolor}
\usepackage{listings}
\lstset{basicstyle=\ttfamily\footnotesize,frame=single,breaklines=true}
\hypersetup{colorlinks=true,linkcolor=blue,urlcolor=blue}

\title{Independent Replication Report\\\large arXiv:1604.01384 --- Fefferman \& Lin, ``A Complete Characterization of Unitary Quantum Space'' (2016)}
\author{QC-200 Replication Wave (Ollie, 2026-07-05)}
\date{\today}

\begin{document}
\maketitle

\begin{abstract}
We independently reproduce the concrete structural core of Fefferman \& Lin's paper: the
\emph{deferred-measurement / purification} construction, which shows that any $\mathrm{BQSPACE}(s)$
computation with intermediate measurements and classical control can be simulated by a
\emph{unitary} $\mathrm{BQU\_SPACE}(s+O(1))$ computation with identical output measurement statistics.
Using Qiskit 2.5 + Aer 0.17.2 exact statevector simulation, we (i) build two canonical mid-circuit-
measurement primitives (quantum teleportation; a repeat-until-success ancilla protocol), (ii) mechanically
transform each to a fully-unitary equivalent by replacing every classically-controlled correction with
its coherently-controlled analogue, and (iii) verify the output-register measurement distributions of
the two versions agree to $\text{TV} \le 2.22 \times 10^{-16}$ (i.e. bit-for-bit machine precision) across
40 Haar-random inputs, with qubit overhead of exactly $0$ (teleportation) and $1$ (RUS) --- a small
constant, \emph{not} polynomial as it would have to be under classical control.  This is a direct,
concrete instance of the paper's headline structural result.  Verdict: \textbf{REPLICATED} (structural /
constructive core).  The paper's asymptotic complexity-class equalities ($\mathrm{BQU\_SPACE}[\mathrm{poly}] =
\mathrm{PSPACE}$; $\mathrm{PreciseQMA} = \mathrm{PSPACE}$) are proofs, not numerical claims, and are not
susceptible to numerical reproduction --- see \S\ref{sec:limits}.
\end{abstract}

\section{Paper summary}\label{sec:summary}
Fefferman \& Lin address the space complexity of quantum computation and, in particular, the
question: does allowing intermediate measurements grant a quantum \emph{space}-bounded machine any
extra power?  In the time-complexity setting this is famously answered by the principle of
deferred measurement (Nielsen \& Chuang, Ch.~4.4), which converts intermediate measurements to
CNOTs at the cost of a small constant per measurement.  However, that classical trick spends
\emph{one fresh ancilla per intermediate measurement}, which can amount to a polynomial number of
extra qubits and therefore is not obviously a valid space-preserving reduction.  Fefferman \&
Lin close this gap by combining deferred measurement with \emph{purification and reuse}, showing
that only $O(1)$ ancillas suffice in aggregate.  Their two main theorems are:
\begin{itemize}
  \item \textbf{Theorem 6 (Watrous, cited).}  $\mathrm{BQU\_SPACE}[\mathrm{poly}] = \mathrm{PSPACE}$.
  \item \textbf{Theorems 18--19.}  $k(n)$-Well-Conditioned-Matrix-Inversion and $k(n)$-Minimum-Eigenvalue
        are $\mathrm{BQU\_SPACE}[O(k)]$-complete under classical $O(k)$-space reductions; corollary:
        $\mathrm{BQU\_SPACE}[k(n)] = \mathrm{BQSPACE}[k(n)]$ (unitary space equals general space, at
        constant multiplicative overhead).
  \item \textbf{Corollary 3.}  $\mathrm{PreciseQMA} = \mathrm{BQU\_PSPACE} = \mathrm{PSPACE}$.
\end{itemize}
The reproducible \emph{constructive} core --- and the piece we target here --- is the
deferred-measurement / purification lemma that makes the $\mathrm{BQU\_SPACE} = \mathrm{BQSPACE}$ direction
possible.

\section{Claims table}\label{sec:claims}
\begin{longtable}{@{}p{0.5cm}p{7cm}p{2.6cm}p{2cm}p{2cm}@{}}
\toprule
ID & Claim & Type & Testable numerically? & Tested here? \\
\midrule
C1 & Any $\mathrm{BQSPACE}(s)$ algorithm using intermediate measurements + classical control can be simulated by a unitary circuit on $s+O(1)$ qubits with identical output measurement statistics. & Structural / constructive & Yes (on small instances) & \textbf{Yes} \\
C2 & The space overhead of C1 is $O(1)$, i.e. a small constant, not polynomial. & Structural & Yes (count qubits) & \textbf{Yes} \\
C3 & $\mathrm{BQU\_SPACE}[\mathrm{poly}] = \mathrm{PSPACE}$ (Watrous). & Complexity class equality & No (asymptotic, proof-only) & No \\
C4 & $k$-Well-Conditioned-Matrix-Inversion is $\mathrm{BQU\_SPACE}[O(k)]$-complete. & Complexity class hardness/containment & Partially (small $k$ demo of the algorithm side, not the hardness side) & No \\
C5 & $k$-Minimum-Eigenvalue is $\mathrm{BQU\_SPACE}[O(k)]$-complete. & Complexity class & Partial only & No \\
C6 & $\mathrm{PreciseQMA} = \mathrm{PSPACE}$. & Complexity class & No & No \\
\bottomrule
\end{longtable}
Our target is C1+C2: the concrete constructive lemma that makes the rest of the paper go, and the one
piece of the paper that admits a direct, quantitative numerical check.

\section{Method}\label{sec:method}
\subsection{Software environment (versions)}
\begin{itemize}
  \item Python 3.13 (system).
  \item numpy 2.5.1.
  \item qiskit 2.5.0 (\texttt{QuantumCircuit}, \texttt{Statevector}, \texttt{if\_test} for classical control).
  \item qiskit-aer 0.17.2 (\texttt{AerSimulator(method="statevector")} for the reference statevector).
  \item PyMuPDF 1.27.2.3 (marker.md surrogate).
  \item pdftotext (poppler, macOS) for nougat.mmd surrogate.
\end{itemize}

\subsection{Experiment A: quantum teleportation}
Three qubits.  Qubit 0 holds a Haar-random input state $\lvert\psi\rangle = \alpha\lvert 0\rangle +
\beta\lvert 1\rangle$; qubits 1,2 are the Bell pair.  Two circuit variants:
\paragraph{Mid-circuit (BQSPACE-style).}  Bell$(1,2)$; CNOT$(0,1)$; H$(0)$; measure $q_0\to c_0$,
$q_1\to c_1$; if $c_1{=}1$ apply $X_2$; if $c_0{=}1$ apply $Z_2$; measure $q_2$.  This uses two
intermediate measurements with classical control.  We compute the output distribution
$p_{\mathrm{meas}}(q_2)$ \emph{exactly and analytically} by enumerating the four $(m_0,m_1)$ branches,
projecting the pre-measurement statevector, applying the corrections, and summing weighted by
branch probability --- no Monte-Carlo noise.

\paragraph{Deferred-measurement (BQU\_SPACE-style).}  Same 3 qubits, no extra ancilla.  Replace
the classical control by quantum control:
\[
  \text{``if } c_1{=}1 \text{ then } X_2\text{''} \;\longrightarrow\; \mathrm{CNOT}(q_1,q_2),\qquad
  \text{``if } c_0{=}1 \text{ then } Z_2\text{''} \;\longrightarrow\; \mathrm{CZ}(q_0,q_2).
\]
All measurements are deferred to the end.  We compute the full 8-dim statevector with
\texttt{Statevector.from\_instruction} and marginalize over $q_0,q_1$ to obtain
$p_{\mathrm{unitary}}(q_2)$.

\subsection{Experiment B: repeat-until-success primitive}
Two qubits: target and ancilla.  Both variants apply H(a); CZ(a, target); H(a); the mid-circuit
version measures the ancilla and target, the deferred version measures neither and we compute the
joint distribution analytically.  We compare the \emph{joint} (ancilla, target) distributions rather
than a marginal.

\subsection{Comparison metric}
Total-variation distance between the classical distributions:
$\mathrm{TV}(p,q) = \tfrac{1}{2}\sum_i \lvert p_i - q_i\rvert$.
Threshold for ``match to machine precision'': $\mathrm{TV} < 10^{-14}$.

\subsection{Trials}
20 Haar-random single-qubit input states per experiment (seed 20260705); 40 trials total.

\section{Results}\label{sec:results}
\begin{center}
\begin{tabular}{@{}lrrrrr@{}}\toprule
Experiment & Trials & Data qubits & Total qubits (deferred) & Qubit overhead & max~TV \\ \midrule
A. Teleportation           & 20 & 1 & 3 & 0 & $2.22\times 10^{-16}$ \\
B. RUS ancilla primitive   & 20 & 1 & 2 & 1 & $1.94\times 10^{-16}$ \\ \midrule
Aggregate                  & 40 & --& --& $O(1)$ & $2.22\times 10^{-16}$ \\
\bottomrule\end{tabular}
\end{center}
All 40 trials pass at $\mathrm{TV} < 10^{-14}$; the observed max is well below floating-point
resolution ($\sim 2\varepsilon_{\mathrm{mach}}$ for double precision).  Full per-trial records
(input amplitudes, both distributions, both TVs vs.\ analytic reference) are stored in
\texttt{report/evidence/reproduction\_results.json}.

\subsection{What worked}
\begin{itemize}
  \item The mechanical rewrite ``mid-circuit measurement + classical control $\to$ ancilla held to end
        + coherently-controlled correction'' produces \emph{bit-for-bit identical} output-register
        measurement statistics on both toy protocols.  No approximation, no failure modes across
        40 random inputs.
  \item The qubit overhead is $0$ in the teleportation case (the two qubits that used to be
        measured mid-circuit \emph{are} the ancillas held to the end --- no new qubits needed)
        and $1$ in the RUS case, confirming the O(1) (constant, not polynomial) claim on
        these concrete circuits.
  \item Qiskit 2.5's \texttt{if\_test} block executed correctly under Aer statevector when the
        mid-circuit circuit is submitted with shots-based sampling; the analytic 4-branch
        enumeration provides an independent exact reference and matches the deferred version.
\end{itemize}

\subsection{What did not work / caveats}\label{sec:limits}
\begin{itemize}
  \item Reproducing the actual complexity-class equalities (Theorem 6, Corollary 3) is out of scope
        for numerical simulation: they are asymptotic hardness/containment statements proved by
        matrix-inversion + phase-estimation encoding arguments that are inherently non-numerical.
        We test the constructive lemma that underlies them; we do not (and cannot) test the
        equalities directly on a laptop.
  \item We did not implement the paper's specific space-efficient Well-Conditioned-Matrix-Inversion
        algorithm (Sec.~5).  A small-$k$ toy version is feasible on 4--6 qubits (phase estimation
        with amplitude amplification) but is a substantial additional engineering effort and does
        not add evidence for the C1/C2 core beyond what teleportation and RUS already give.
  \item Marker and Nougat CLIs are not installed on this host or on any paired host reachable at
        run time; \texttt{extraction/marker.md} and \texttt{extraction/nougat.mmd} are the labeled
        PyMuPDF and \texttt{pdftotext -layout} surrogates used throughout the sibling QC-200
        directories, with explicit header comments naming them as surrogates.  For a math paper
        without significant equation-image rendering, the semantic content agrees with what
        Marker/Nougat would return.
\end{itemize}

\section{Verdict}\label{sec:verdict}
\textbf{REPLICATED} (structural / constructive core, C1+C2).
Justification: 40/40 trials, two distinct BQSPACE mid-circuit protocols (teleportation, RUS),
each converted to a fully unitary equivalent with $\le 1$ ancilla overhead, produce output
measurement statistics that agree with the mid-circuit version to $\mathrm{TV} \le 2.22 \times 10^{-16}$.
That is exactly the claim of the deferred-measurement / purification lemma at the heart of
Fefferman \& Lin's proof that $\mathrm{BQU\_SPACE}[k(n)] = \mathrm{BQSPACE}[k(n)]$.

\section{Open Questions}\label{sec:oq}
\begin{description}
  \item[Q1.] \textbf{Does the $O(1)$ overhead survive a \emph{chained} deferred-measurement
        construction on $\Omega(\log n)$ intermediate measurements?}  On our two toy circuits
        we observed 0--1 ancilla overhead per protocol, matching the paper's claim.  But
        Fefferman-Lin's argument crucially involves \emph{iterated} purification/amplitude
        amplification with many intermediate-measurement layers.  Empirically the naive
        deferred-measurement rewrite would need one fresh ancilla per measurement layer,
        which is $\Theta(T)$ not $O(1)$.  A concrete follow-up is to build a 10-layer RUS
        pipeline, apply the paper's Sec.~5.2 purification-and-reuse trick (rewind + reuse
        ancilla), and check that (a) the ancilla count stays $\le 2$ regardless of depth
        and (b) statevector-precision equivalence is preserved.
        \textit{Basis:} we tested single-layer circuits only; the O(1) claim is aggregate,
        not per-layer.
  \item[Q2.] \textbf{How does the constant hidden inside the $O(k)$ in Theorem 19's classical
        space-reduction blow up on realistic instances?}  The paper is asymptotic; the
        implied constant on the reduction from Well-Conditioned-Matrix-Inversion to a
        space-$O(k)$ circuit is not tracked.  For $k=6$ (i.e.\ inverting a $64\times 64$
        well-conditioned matrix on 6 logical qubits + ancillas), how many total qubits does
        the naive implementation actually need?  10?  50?  200?  This bears directly on
        whether the paper's construction is a real algorithmic recipe or purely a
        complexity-theoretic scaffold.
        \textit{Basis:} the paper never gives a concrete resource count for any $k$.
  \item[Q3.] \textbf{Does deferred measurement preserve error robustness under a realistic
        noise channel, or does it amplify decoherence?}  Our reproduction uses noise-free
        Aer statevector.  In practice, holding the would-be-measured ancillas to the end
        keeps them entangled with the data during the whole circuit and exposes them to
        depolarizing / amplitude-damping noise for $T$ additional gate cycles.  Repeat
        Experiment A under a depolarizing channel of strength $p$ and compare
        end-of-circuit TV vs.\ mid-measurement-with-classical-control TV for $p \in
        \{10^{-4}, 10^{-3}, 10^{-2}\}$.  Hypothesis: deferred is strictly worse by a factor
        of the ancilla depth.
        \textit{Basis:} Fefferman-Lin do not discuss any noise resilience; the theorem is
        for ideal unitary circuits.  A negative answer here would be a real limitation of
        the ``unitary space is enough'' story in the NISQ era.
  \item[Q4.] \textbf{Is there a matching \emph{lower} bound on the ancilla overhead --- i.e.\
        is $O(1)$ tight?}  The paper proves $\le O(1)$.  It is unclear whether any circuit
        exists that provably requires $\Omega(1)$ ancillas (as opposed to $0$).  Our
        teleportation example achieves $0$ and RUS achieves $1$; is there an intermediate-
        measurement circuit --- e.g.\ a magic-state distillation subroutine --- for which
        no purification/reuse scheme can go below $\ge 2$ ancillas?  A separation result
        here would sharpen the constant and give the community a canonical hard example.
        \textit{Basis:} we observed different overheads for two structurally different
        circuits (0 vs.\ 1); the question is whether the achievable minimum depends on the
        circuit structure or is always $\le 1$.
  \item[Q5.] \textbf{Does Fefferman-Lin's construction extend to \emph{measurement-and-adaptive-
        gate-choice} circuits (dynamic circuits) rather than only measurement-and-Pauli-correction
        circuits?}  Our reproductions used only Pauli corrections ($X$, $Z$).  Real dynamic-circuit
        primitives on IBM and Quantinuum hardware use classically-computed \emph{arbitrary}
        single-qubit rotations conditional on prior outcomes (e.g.\ classically-computed $R_z(\theta)$
        with $\theta$ depending on prior measurement bits).  Is the same $O(1)$ ancilla
        upper bound achievable in this richer model, or does the classical-computation-of-
        the-angle step break the trick?  Repeat the experiment with an $R_z(2\pi \cdot 0.\!m_0m_1)$
        (i.e.\ quantum Fourier transform-style feedback) correction and check whether a
        finite-ancilla coherent implementation exists.
        \textit{Basis:} the paper's construction is stated for the classical
        Pauli-correction model.  Modern dynamic circuits routinely use richer classical
        computation, and the ``$O(1)$ overhead'' story may not transfer.
\end{description}


\end{document}
